\subsection{Inverse Kinematics Approaches for Landmark-Driven Hand Pose Recovery}
\label{subsec:sota_ik}

Given a set of three-dimensional hand landmarks --- such as the 21
keypoints produced by MediaPipe Hands~\cite{zhang2020mediapipehands}
--- recovering the MANO joint rotations~$\boldsymbol{\theta}$ is an
inverse kinematics (IK) problem: find the joint angles whose
forward-kinematics map reproduces the observed positions.
The formulation departs from the classical robotics setting in one
important respect: the 21-landmark case is \emph{fully observed}.
Every vertex of the kinematic tree has a corresponding detected
position; no intermediate joints are latent.
This distinction has direct consequences for which solver families are
appropriate.
The four main approaches surveyed below --- iterative numerical methods,
optimisation-based model fitting, neural regression, and analytical
closed-form solvers --- are evaluated both on their own reported
performance and against the deployment constraints of this work.



\subsubsection{Iterative Numerical Solvers}\label{subsubsec:sota_ik_iterative}

Iterative IK solvers were developed to handle the under-constrained
case: a manipulator arm whose end-effectors are specified but whose
intermediate joints are free.
FABRIK~\cite{aristidou2011fabrik} alternates forward and backward passes
along the kinematic chain, repositioning each joint toward the next
target position until all end-effectors lie within a prescribed
tolerance.
Cyclic-coordinate descent sweeps the chain joint by joint,
each time minimising the displacement at the current degree of freedom
while holding the rest fixed.
Jacobian-based methods linearise the forward kinematics map and invert
it via the pseudo-inverse or the Levenberg--Marquardt damped
least-squares variant~\cite{aristidou2018iksurvey}.
Each of these families has been applied to articulated hand models:
Heap and Hogg~\cite{heap1996dynamic} tracked a deformable hand model
from a single video camera as early as 1996, recovering
six-degree-of-freedom rigid pose together with PCA-based shape
deformation through iterative weighted-least-squares fitting.

The fully-observed 21-landmark case exposes a structural mismatch,
however.
Iteration is necessary when the problem is under-constrained; when every
joint position is already known, the kinematic chain is fully
determined and no iteration is required to infer missing data.
Applied to the 21-landmark setting, an iterative solver converges to a
solution that the input already encodes, incurring convergence-dependent
overhead without recovering any additional information.
For a real-time pipeline with a tight per-frame latency budget on CPU,
this is a gratuitous cost.



\subsubsection{Optimisation-Based Model Fitting}\label{subsubsec:sota_ik_optim}

A second approach formulates IK as minimising a cost function over MANO
parameters: find $\boldsymbol{\theta}$ (and optionally
$\boldsymbol{\beta}$) whose forward-kinematics output best matches the
detected keypoints, subject to anatomical joint-limit constraints.
The objective may include a 3D positional term, a 2D reprojection term,
or both, and the minimisation is carried out by a gradient-based solver
such as L-BFGS.

Drosakis and Argyros~\cite{drosakis2023mano} demonstrate this strategy
by fitting MANO parameters to OpenPose 2D hand keypoints using L-BFGS
under anatomical joint-limit constraints and shape regularisation,
reporting accuracy competitive with learning-based methods and superior
to other optimisation-based methods on the Dexter+Object and EgoDexter
benchmarks.
The approach is entirely data-free and produces MANO output that
respects anatomical priors by construction, two properties difficult
to achieve with purely learned methods.
Earlier work by Panteleris et al.~\cite{panteleris2018backrgb} showed
that optimisation over a kinematic hand model from RGB input could
produce usable pose estimates without depth data, establishing
the viability of the model-fitting paradigm for monocular settings.

The principal limitation of gradient-based fitting is that per-frame
solve time is data-dependent: the number of iterations required varies
with pose complexity, the initialisation quality, and the curvature of
the loss landscape.
On CPU hardware, where evaluating the MANO forward pass is
considerably slower than on GPU, this makes the solver unsuitable for
interactive-rate applications where a consistent per-frame budget must
be met.



\subsubsection{Neural Regression}\label{subsubsec:sota_ik_neural}

Neural IK replaces the iterative or gradient loop with a single
feed-forward pass, bounding per-frame inference cost regardless of pose
complexity.

Zhou et al.~\cite{zhou2020monocular} introduced IKNet as part of the
Minimal-Hand pipeline: a seven-layer fully-connected network trained on
motion-capture data that maps root-relative 3D joint positions directly
to per-joint unit-quaternion rotations.
IKNet adds approximately \SI{0.9}{\ms} of inference overhead on an
NVIDIA GTX\,1080\,Ti, and the full Minimal-Hand pipeline
(DetNet\,+\,IKNet) exceeds 100\,fps on that
hardware~\cite{zhou2020monocular}.
The network is compact and fast; inference cost is strictly bounded by
the forward-pass time regardless of the input pose.
Published performance figures are on GPU hardware, and the network
encodes a pose prior from training data, introducing a data dependency
absent from purely geometric methods.

Li et al.~\cite{li2021hybrik} refined the neural approach by identifying
which component of the rotation actually requires learning.
Each per-joint rotation can be decomposed into a \emph{swing} --- the
minimal rotation aligning the canonical bone direction with the
observed bone direction, uniquely and analytically determined from the
3D landmarks --- and a \emph{twist}, a rotation about the bone's own
axis that the swing leaves underdetermined.
HybrIK solves the swing in closed form from the landmarks and estimates
the twist from image features via a lightweight network, achieving
state-of-the-art body-mesh recovery on 3DPW at CVPR\,2021.
The architecture makes explicit what must be learned (the twist) and
what need not be (the swing), substantially reducing the capacity
required of the network.
Li et al.\ subsequently extended the same decomposition to whole-body
recovery, including articulated hands, in
HybrIK-X~\cite{li2023hybrikx}.
Both models are evaluated exclusively on GPU hardware in their published
benchmarks; no interactive-rate deployment on CPU is reported for either.



\subsubsection{Analytical Closed-Form Solvers}\label{subsubsec:sota_ik_analytical}

The fourth family removes the neural component from the IK path
entirely, relying instead on the geometric structure of the
fully-observed landmark set.

The AIK algorithm shares the swing–twist decomposition that HybrIK
also employs (originally due to Baerlocher and Bouloch) but diverges
at the twist.
Global root orientation is recovered as the rotation that best aligns
the five canonical wrist-to-MCP vectors with their observed
counterparts, solved in closed form via singular-value decomposition
of the cross-covariance matrix following
Arun et al.~\cite{arun1987leastsquares}.
Each subsequent joint rotation is then resolved as a pure swing along
the kinematic chain: the minimal rotation mapping the parent-relative
canonical bone direction onto the observed direction.
The twist is set to zero for all finger joints.

This is not a numerical shortcut.
Axial rotation about a finger metacarpal or phalanx is
physiologically negligible in natural hand motion; the dominant degrees
of freedom at each finger joint are flexion--extension and, at the
metacarpophalangeal joints, abduction--adduction, both of which are
fully captured by the swing~(\autoref{subsec:aik}).
The discarded twist degree of freedom carries minimal perceptual or
kinematic weight in practice.

With the iterative loop, the gradient solver, and the neural network
all removed, the entire landmark-to-rotation path reduces to a single
traversal of the kinematic tree.
Per-frame cost is deterministic, scales only with the number of joints,
and carries no dependency on training data, convergence tolerance, or
GPU hardware.



\subsubsection{Summary and Solver Selection}\label{subsubsec:sota_ik_summary}

Iterative solvers are data-free but impose convergence overhead that is
structurally unnecessary when the input is fully observed.
Optimisation-based fitting produces anatomically regularised output but
its unbounded solve time disqualifies it from interactive CPU
deployment.
Neural regression, and in particular the HybrIK-X decomposition,
achieves the highest reconstruction accuracy, recovering the twist
component that geometry alone cannot determine; the cost is a GPU
backbone for the twist network that places it outside the deployment
target.
AIK is the only approach satisfying all three constraints
simultaneously: it is data-free, executes in deterministic time on
CPU, and its per-frame cost is bounded by construction.
HybrIK-X therefore defines the accuracy ceiling the present system does
not reach; AIK defines the latency and hardware floor it operates within.
